\documentclass[11pt,letterpaper]{article}
\PassOptionsToPackage{hyphens}{url}
\usepackage[margin=.83in,top=.8in,bottom=.8in,headheight=15pt]{geometry}
\usepackage[T1]{fontenc}
\usepackage{lmodern,amsmath,amssymb,mathtools,bm,booktabs,array,graphicx,tikz}
\usetikzlibrary{arrows.meta,positioning,calc,decorations.pathreplacing,patterns}
\definecolor{ink}{HTML}{17344D}\definecolor{accent}{HTML}{087E8B}\definecolor{pale}{HTML}{EDF5F5}\definecolor{gold}{HTML}{AD6B19}
\providecommand{\dd}{\,\mathrm{d}}\providecommand{\R}{\mathbb{R}}
\providecommand{\E}{\mathbb{E}}\providecommand{\Var}{\operatorname{Var}}
\providecommand{\Cov}{\operatorname{Cov}}\providecommand{\diag}{\operatorname{diag}}
\providecommand{\norm}[1]{\left\lVert#1\right\rVert}
\providecommand{\abs}[1]{\left|#1\right|}
\usepackage{microtype,hyperref,fancyhdr,enumitem,amsthm,sectsty,needspace}
\allsectionsfont{\raggedright}
\hypersetup{colorlinks=true,linkcolor=ink,urlcolor=accent,pdfauthor={Zan Ahmad},pdftitle={Local convolutional segmentation and geometric error / Review}}
\setlength{\parindent}{0pt}\setlength{\parskip}{6pt}\linespread{1.07}
\setlist{itemsep=3pt,topsep=4pt,leftmargin=*}
\emergencystretch=2em\clubpenalty=10000\widowpenalty=10000
\pagestyle{fancy}\fancyhf{}\fancyhead[L]{\small Computational Cardiology}\fancyhead[R]{\small Lecture 12}
\fancyfoot[L]{\footnotesize Zan Ahmad}\fancyfoot[R]{\footnotesize\thepage}\renewcommand{\headrulewidth}{.3pt}
\newtheorem{proposition}{Proposition}\newtheorem{theorem}{Theorem}\newtheorem{lemma}{Lemma}
\theoremstyle{definition}\newtheorem{definition}{Definition}\newtheorem{example}{Example}
\numberwithin{equation}{section}
\begin{document}
\section*{Local convolutional segmentation and geometric error --- review sheet}
Use this after attempting the worked examples and practice. Explain each claim in words before calculating.

\section*{Key ideas}
\begin{itemize}
\item A shared 3$\times$3 layer learns ten parameters including bias.
\item Loss gradients can be checked numerically before trusting a training run.
\item The split unit should match the intended unit of generalization.
\item Overlap, boundary distance, area and connectivity describe different failures.
\item Synthetic improvement is evidence about the declared generator, not clinical validation.
\end{itemize}
\section*{Common errors and repairs}
\textbf{Call the small scaffold U-Net.} It has one trained convolution and sigmoid; the U-Net reading is an architecture comparison.

\textbf{Tune a threshold after inspecting test masks.} Choose with training/validation information, freeze it, then evaluate the test protocol.

\textbf{Treat a missing predicted object as zero boundary error.} Mark the nonempty/empty distance unavailable and report the failure.

\textbf{Accept high Dice as proof that topology is preserved.} Check the feature relevant to the downstream question, such as connected components.

\section*{Retrieve without your notes}
\begin{itemize}
\item Why can area error be zero with an incorrect boundary?
\item What does a central finite-difference check compare?
\item How does this assignment fit the course policy?
\end{itemize}
\section*{Readiness check}
\begin{itemize}
\item Distinguish a fixed intensity baseline from a trained shared-convolution segmenter.
\item Trace a ten-parameter model from local image patches through a loss and gradient update.
\item Evaluate complete held-out images using overlap, physical boundary distances and connectivity.
\item Combine geometry provenance and a sensitivity counterexample in a short Assignment 3 report.
\end{itemize}
\end{document}
